\documentclass[openany,fontset=fandol]{ctexbook}
\usepackage{amsmath,amssymb}
\usepackage[unicode,hidelinks]{hyperref}
\setlength{\emergencystretch}{2em}
\newcommand{\Sym}{\operatorname{Sym}}
\newcommand{\gr}{\operatorname{gr}}
\begin{document}
\title{代数几何基础第五卷预备笔记\\划去稿：命题8与引理8.1}
\author{据英译编辑稿续译，参校法文扫描}
\date{2026年11月20日}
\maketitle
\chapter*{范围与接续}
本批从英译第69页命题8开始，连续译至引理8.1后的证明段末；对读法文第105页。下一起点为同页命题9。下文仍属作者划去稿，未并入第十七版第五卷增补。正文保留原式及疑义，后列数学核查。
\chapter*{作者划去稿（续）}
\noindent\textbf{命题8。}
增广$A\to A/I$为普通二次型的充要条件，是满足以下各条：
\begin{enumerate}
\item[$1^\circ$] 各$\gr_I^n(A)$都是有限型投射$A/I$-模。
\item[$2^\circ$] 若以$d$表示$\gr_I^1(A)=I/I^2$在$p\in\operatorname{Spec}(A/I)$处的秩，则$\gr_I^n(A)$在同一点处的秩等于（或者只需大于或等于）
\[
\binom{n+d}{d}-\binom{n-2+d}{d}.
\]
\item[$3^\circ$] （$1^\circ$及$2^\circ$已蕴含$K=\ker\theta^2$是可逆模，在$\Sym^2(I/I^2)$中局部为直和因子。）$\Sym^2(I/I^2)$的子模$K=\ker\theta^2$是非退化的。
\end{enumerate}

借助下述引理，证明是直接的。

\noindent\textbf{引理8.1。}
设$B$是一个环，$M$是一个有限型投射$B$-模，$\phi\in\Sym^2(M)$是一个“非退化”元素。那么，乘以$\phi$的同态
\[
\Sym^{[\text{次数待核}]}(M)\longrightarrow
\Sym^{[\text{次数待核}]}(M)
\]
是到某个直和因子子模的同构。因此，若$M$的秩为$d$，其余核是秩为
\[
\binom{n+d}{d}-\binom{n-2+d}{d}
\]
的投射模。\footnote{英译者已注明两处上标难辨，猜测可能涉及$n$和$n+2$。余核首项的上标英译印为$n+3$，法文第105页清楚为$n+d$，故本处据法文；此项仍保留原来的下标$d$。}

我们立刻归结到$B$是一个域的谱的情形，这时结论是平凡的。\footnote{法文也字面写$B$是一个域的谱，但$B$此前是环。语境中的归约是逐点取剩余域（其谱只有一点），并非将$B$等同于其谱。}（实际上，“非退化”只用来保证各个$\phi(x)$，$x\in\operatorname{Spec}(B)$，都不为零。用任意整数$k\ge1$代替$2$，而不提非退化，或许能更合理地表述本引理。）
\chapter*{译校说明（非原文）}
\paragraph{命题8的秩障碍。}
法文明确补出$d$为一次分量的秩，且两项的二项式下标均为$d$，也保留“或只需$\ge$”。不能只将这些缺陷归于英译。令$n=2$，原式化为$\binom{d+2}{d}-1=d(d+3)/2$；而$\theta^2$为满射，故二次分量秩至多$\operatorname{rk}\Sym^2(I/I^2)=d(d+1)/2$。当$d\ge1$，前者比后者多$d$，连下界也不可能。由此不能沿用原式推出$3^\circ$括号中的结论。

\paragraph{引理的可复原内容。}
法文两次数上标有叠写，英译者也标明难辨。按目标处于$n$次来重建，乘二次式的映射应为
$\Sym^{n-2}(M)\to\Sym^n(M)$（$n<2$时源为零）；也可把源次数记为$n$，则目标为$n+2$。这不是辨读出法文最终字样。若$\phi(x)$在每个剩余域上非零，纤维的对称代数为多项式环，乘$\phi(x)$为单射。将有限投射模局部平凡化后，源秩阶的某个子式在每一点附近可逆，故原映射局部裂单，余核有限局部自由，从而是有限投射模。固定正秩$d$时，所得余核秩为
\[
h_n=\binom{n+d-1}{d-1}-\binom{n-2+d-1}{d-1},
\]
其中负次数的对称幂按零计。若改用次数$k$，同一论证只要求每个纤维中的$k$次式非零。

\paragraph{即使改正计数，下界也不能独自保证有二次关系。}
设$A=k[T_1,\dots,T_d]$，$I=(T_1,\dots,T_d)$，$d\ge1$。则$\gr_I(A)=\Sym_k(I/I^2)$；每个分量均有限维且维数至少为$h_n$，但$K=0$，并非可逆模。这反驳“$1^\circ$和下界已推出$K$可逆”的说法，不反驳另行假定$K$为可逆且非退化后的判别。后者在有限投射及各纤维中$K$的局部生成元非零的条件下，先由上述引理给出$\Sym(I/I^2)/(K)$各次秩为$h_n$，再与其到$\gr_I(A)$的各次满射比较；此时$\gr_I^n(A)$的秩下界$h_n$足以使满射成为同构。这些是数学重建，不是作者已经写出的替换命题。
\end{document}
